\section{Analysis}
\label{sec:analysis}

This section consolidates cross-cutting observations that did not fit
naturally inside any single capability section.

\subsection{Headline-WER Failure Mode: FLEURS-en}

Among the English benchmarks in Table~\ref{tab:asr_en}, FLEURS-en
stands out at $30.47$\% WER -- substantially higher than CommonVoice-en
($10.12$\%) and the LibriSpeech sets. FLEURS uses prompted-read speech
of factual sentences across many languages; the gap is consistent with
two well-known phenomena
% TODO: verify the breakdown by qualitatively inspecting a sample of
% FLEURS-en errors before keeping these explanations in the final
% version.
\begin{enumerate}
  \item the English slice of FLEURS contains a high density of named
        entities (toponyms, agencies) that are rare in our training
        mix, which would also explain why the hotword path closes most
        of this gap (see CV17-en hotword results,
        Section~\ref{sec:hotwords});
  \item FLEURS reference texts mix punctuation conventions that our
        normalizer does not fully canonicalise.
\end{enumerate}
The next training iteration will (i)~increase coverage of read,
factual English to address the first effect and (ii)~align the FLEURS
text-normalization protocol with the Whisper normalizer used by recent
audio-LLM reports.

\subsection{Hotword Curve: Why Five}
\label{subsec:why-five}

Section~\ref{sec:hotwords} showed that the optimal in-prompt hotword
count is around five. We attribute this to two competing effects:

\begin{itemize}
  \item \emph{Recall effect.} Each true-positive hotword in the prompt
        is, in expectation, an additional rare term the model can
        ``snap'' to. This effect saturates quickly because the typical
        utterance contains at most one or two rare entities.
  \item \emph{False-trigger effect.} Each distractor hotword is a
        candidate insertion / substitution that the model must
        \emph{not} commit to. The risk grows roughly linearly with
        list length.
\end{itemize}

Below five, the recall effect dominates; above five, the false-trigger
effect catches up.
% TODO: support this story with an empirical decomposition into
% insertion vs. substitution errors per k; the data is in the eval
% logs but is not yet plotted.

\subsection{Libri2Mix vs Internal TS-ASR}

The gap between $14.29$\% MER on the internal target-speaker test and
$65.83$\% WER on Libri2Mix is the most important open issue in this
release. Section~\ref{sec:tsasr} narrows the suspect list down to a
distributional mismatch in overlap fraction, SNR distribution and the
training-time RIR application rate; the planned ablation -- one
training epoch with overlap drawn uniformly from $[0.5, 1.0]$, plus
inference with RIR disabled -- will be the cheapest way to attribute
the gap before the next training iteration.

\subsection{Inherited Failure Modes}

LLM-based ASR is known to inherit characteristic LLM failure modes
\cite{radford2023whisper, chu2023qwenaudio}. We have observed (and the
next release will measure):

\begin{itemize}
  \item Repetition loops on long silences or stationary noise. The
        noise-only negative branch of the TS-ASR pipeline
        (Section~\ref{subsec:tsasr-synth}) explicitly trains the model
        to terminate on silence under the TS prompt; we have not yet
        confirmed whether the same protection transfers to the
        plain-ASR prompt without further work
        (Section~\ref{sec:noise}).
  \item Language drift on Mandarin--English code-switching, where short
        English phrases inside a Mandarin utterance can be ``corrected''
        into Mandarin transcriptions of similar-sounding words;
  \item Occasional over-bias on supplied hotwords when the audio does
        not actually contain the term (see
        Section~\ref{subsec:why-five}).
\end{itemize}
% TODO: include a small qualitative table of one example per failure
% mode (sanitised), with audio identifiers traceable to the released
% eval pack.
